% GENERATED FILE. Edit canonical lessons, then run npm run generate:books.
% canonical-source-hash: fnv1a64:47eac9c2bc7dd799
% canonical-lessons: TA-C474-arikuri, TA-C474-kumattal, TA-C474-patapatappu, TA-C474-vikku, TA-C474-uruttu-chafe

\chapter{Signs of Illness}
\label{ch:ta-signs-of-illness}
\hlchaptermodality{474}

\begin{chapteropening}
\textbf{By the end of this chapter:} \emph{I can describe symptoms such as nausea, palpitations and hiccups.}

Complete the last lesson of chapter 474: I can describe symptoms such as nausea, palpitations and hiccups.
\end{chapteropening}

\section[\texorpdfstring{\ta{அறிகுறி} (aṟikuṟi)}{அறிகுறி (aṟikuṟi)}]{\ta{அறிகுறி} (aṟikuṟi) — a symptom; a sign}

\label{lesson:TA-C474-arikuri}



\begin{quote}
Before the new one: say the Tamil for self-confidence, then the Tamil for to admire.
\end{quote}

\subsection*{You'll want to know: \ta{அறிகுறி}}
\textbf{\ta{அறிகுறி}} — \emph{aṟikuṟi} — \textquotedblleft{}a symptom; a sign\textquotedblright{}.

\begin{quote}
\ta{காய்ச்சல்} \ta{இந்த} \ta{நோயின்} \ta{முதல்} \ta{அறிகுறி}.
\end{quote}

\begin{quote}
\emph{kāyccal inta nōyiṉ mutal aṟikuṟi.}
\end{quote}

\begin{quote}
Fever is the first symptom of this illness.
\end{quote}

Plural: \ta{அறிகுறிகள்}, \emph{aṟikuṟikaḷ}.

\ta{அறி} \textquotedblleft{}to know\textquotedblright{} + \ta{குறி} \textquotedblleft{}a mark\textquotedblright{}: a mark that lets you know.

\begin{cousinweb}
\textbf{\ta{அறிகுறி}} (\emph{aṟikuṟi}) is a native compound of \ta{அறி} \emph{aṟi}, \textquotedblleft{}to know\textquotedblright{} (DEDR 314), and \ta{குறி} \emph{kuṟi}, \textquotedblleft{}a mark, a sign\textquotedblright{} (DEDR 1847).
\end{cousinweb}

\begin{grammarlens}[title={What you've built}]
One more word for describing what you lived and felt.
\end{grammarlens}

\subsection*{Your turn}
\begin{itemize}
\raggedright
  \item \emph{Say it:} \emph{aṟikuṟi}
  \item \emph{Say it:} \emph{aṟikuṟi}, once more
  \item \emph{Say it:} \emph{kāyccal inta nōyiṉ mutal aṟikuṟi.} — Fever is the first symptom of this illness.
  \item \emph{Recall:} say the Tamil for self-confidence, then the Tamil for to admire, then say \emph{aṟikuṟi} again
\end{itemize}

\begin{tcolorbox}[breakable,colback=teal!4,colframe=teal!35!black,title={Before you move on}]
What does \ta{அறிகுறி} mean? (\textquotedblleft{}A symptom; a sign\textquotedblright{}.) Say it once more.
\end{tcolorbox}
\section[\texorpdfstring{\ta{குமட்டல்} (kumaṭṭal)}{குமட்டல் (kumaṭṭal)}]{\ta{குமட்டல்} (kumaṭṭal) — nausea}

\label{lesson:TA-C474-kumattal}



\begin{quote}
Before the new one: say the Tamil for to admire, then the Tamil for a symptom.
\end{quote}

\subsection*{You'll want to know: \ta{குமட்டல்}}
\textbf{\ta{குமட்டல்}} — \emph{kumaṭṭal} — \textquotedblleft{}nausea\textquotedblright{}.

\begin{quote}
\ta{பேருந்தில்} \ta{பயணம்} \ta{செய்தபோது} \ta{எனக்குக்} \ta{குமட்டல்} \ta{வந்தது}.
\end{quote}

\begin{quote}
\emph{pēruntil payaṇam ceytapōtu eṉakkuk kumaṭṭal vantatu.}
\end{quote}

\begin{quote}
I felt sick while travelling by bus.
\end{quote}

Dative experiencer: \ta{எனக்குக்} \ta{குமட்டல்} \ta{வந்தது}, I felt nauseous.

\begin{cousinweb}
\textbf{\ta{குமட்டல்}} (\emph{kumaṭṭal}) is native Dravidian, a noun from \ta{குமட்டு} \emph{kumaṭṭu}, \textquotedblleft{}to retch, to feel like vomiting\textquotedblright{} (DEDR 1739).
\end{cousinweb}

\begin{grammarlens}[title={What you've built}]
One more word for describing what you lived and felt.
\end{grammarlens}

\subsection*{Your turn}
\begin{itemize}
\raggedright
  \item \emph{Say it:} \emph{kumaṭṭal}
  \item \emph{Say it:} \emph{kumaṭṭal}, once more
  \item \emph{Say it:} \emph{pēruntil payaṇam ceytapōtu eṉakkuk kumaṭṭal vantatu.} — I felt sick while travelling by bus.
  \item \emph{Recall:} say the Tamil for to admire, then the Tamil for a symptom, then say \emph{kumaṭṭal} again
\end{itemize}

\begin{tcolorbox}[breakable,colback=teal!4,colframe=teal!35!black,title={Before you move on}]
What does \ta{குமட்டல்} mean? (\textquotedblleft{}Nausea\textquotedblright{}.) Say it once more.
\end{tcolorbox}
\section[\texorpdfstring{\ta{படபடப்பு} (paṭapaṭappu)}{படபடப்பு (paṭapaṭappu)}]{\ta{படபடப்பு} (paṭapaṭappu) — palpitations; agitation}

\label{lesson:TA-C474-patapatappu}



\begin{quote}
Before the new one: say the Tamil for a symptom, then the Tamil for nausea.
\end{quote}

\subsection*{You'll want to know: \ta{படபடப்பு}}
\textbf{\ta{படபடப்பு}} — \emph{paṭapaṭappu} — \textquotedblleft{}palpitations; agitation\textquotedblright{}.

\begin{quote}
\ta{மேடையில்} \ta{பேசும்போது} \ta{எனக்குப்} \ta{படபடப்பாக} \ta{இருந்தது}.
\end{quote}

\begin{quote}
\emph{mēṭaiyil pēcumpōtu eṉakkup paṭapaṭappāka iruntatu.}
\end{quote}

\begin{quote}
My heart was racing while I spoke on the stage.
\end{quote}

Of a racing heart, or of nerves before an event.

\begin{cousinweb}
\textbf{\ta{படபடப்பு}} (\emph{paṭapaṭappu}) is native Dravidian (DEDR 3842), the noun of a doubled verb \ta{படபட} \emph{paṭapaṭa}, \textquotedblleft{}to be over-hasty\textquotedblright{}.
\end{cousinweb}

\begin{grammarlens}[title={What you've built}]
One more word for describing what you lived and felt.
\end{grammarlens}

\subsection*{Your turn}
\begin{itemize}
\raggedright
  \item \emph{Say it:} \emph{paṭapaṭappu}
  \item \emph{Say it:} \emph{paṭapaṭappu}, once more
  \item \emph{Say it:} \emph{mēṭaiyil pēcumpōtu eṉakkup paṭapaṭappāka iruntatu.} — My heart was racing while I spoke on the stage.
  \item \emph{Recall:} say the Tamil for a symptom, then the Tamil for nausea, then say \emph{paṭapaṭappu} again
\end{itemize}

\begin{tcolorbox}[breakable,colback=teal!4,colframe=teal!35!black,title={Before you move on}]
What does \ta{படபடப்பு} mean? (\textquotedblleft{}Palpitations; agitation\textquotedblright{}.) Say it once more.
\end{tcolorbox}
\section[\texorpdfstring{\ta{விக்கு} (vikku)}{விக்கு (vikku)}]{\ta{விக்கு} (vikku) — to hiccup}

\label{lesson:TA-C474-vikku}



\begin{quote}
Before the new one: say the Tamil for nausea, then the Tamil for palpitations.
\end{quote}

\subsection*{You'll want to know: \ta{விக்கு}}
\textbf{\ta{விக்கு}} — \emph{vikku} — \textquotedblleft{}to hiccup\textquotedblright{}.

\begin{quote}
\ta{வேகமாகச்} \ta{சாப்பிட்டதால்} \ta{எனக்கு} \ta{விக்கியது}.
\end{quote}

\begin{quote}
\emph{vēkamākac cāppiṭṭatāl eṉakku vikkiyatu.}
\end{quote}

\begin{quote}
I got hiccups because I ate fast.
\end{quote}

Past (it …): \ta{விக்கியது}, \emph{vikkiyatu}.

Dative experiencer, so the past is \ta{விக்கியது}. \ta{விக்கல்}, a hiccup.

\begin{cousinweb}
\textbf{\ta{விக்கு}} (\emph{vikku}) is native Dravidian (DEDR 5383); the entry gives \ta{விக்கல்} \emph{vikkal}, \textquotedblleft{}a hiccup\textquotedblright{}.
\end{cousinweb}

\begin{grammarlens}[title={What you've built}]
One more word for describing what you lived and felt.
\end{grammarlens}

\subsection*{Your turn}
\begin{itemize}
\raggedright
  \item \emph{Say it:} \emph{vikku}
  \item \emph{Say it:} \emph{vikku}, once more
  \item \emph{Say it:} \emph{vēkamākac cāppiṭṭatāl eṉakku vikkiyatu.} — I got hiccups because I ate fast.
  \item \emph{Recall:} say the Tamil for nausea, then the Tamil for palpitations, then say \emph{vikku} again
\end{itemize}

\begin{tcolorbox}[breakable,colback=teal!4,colframe=teal!35!black,title={Before you move on}]
What does \ta{விக்கு} mean? (\textquotedblleft{}To hiccup\textquotedblright{}.) Say it once more.
\end{tcolorbox}
\section[\texorpdfstring{\ta{உறுத்து} (uṟuttu)}{உறுத்து (uṟuttu)}]{\ta{உறுத்து} (uṟuttu) — to chafe; to irritate}

\label{lesson:TA-C474-uruttu-chafe}



\begin{quote}
Before the new one: say the Tamil for palpitations, then the Tamil for to hiccup.
\end{quote}

\subsection*{You'll want to know: \ta{உறுத்து}}
\textbf{\ta{உறுத்து}} — \emph{uṟuttu} — \textquotedblleft{}to chafe; to irritate\textquotedblright{}.

\begin{quote}
\ta{கண்ணில்} \ta{தூசி} \ta{விழுந்து} \ta{உறுத்தியது}.
\end{quote}

\begin{quote}
\emph{kaṇṇil tūci viḻuntu uṟuttiyatu.}
\end{quote}

\begin{quote}
Dust got into my eye and irritated it.
\end{quote}

Past (it …): \ta{உறுத்தியது}, \emph{uṟuttiyatu}.

Of the conscience too: \ta{மனம்} \ta{உறுத்தியது}, my conscience troubled me.

\begin{cousinweb}
\textbf{\ta{உறுத்து}} (\emph{uṟuttu}) is native Dravidian: DEDR 712 lists \ta{உறுத்து} with the sense of pressing hard upon the mind.
\end{cousinweb}

\begin{grammarlens}[title={What you've built}]
One more word for describing what you lived and felt.
\end{grammarlens}

\subsection*{Your turn}
\begin{itemize}
\raggedright
  \item \emph{Say it:} \emph{uṟuttu}
  \item \emph{Say it:} \emph{uṟuttu}, once more
  \item \emph{Say it:} \emph{kaṇṇil tūci viḻuntu uṟuttiyatu.} — Dust got into my eye and irritated it.
  \item \emph{Recall:} say the Tamil for palpitations, then the Tamil for to hiccup, then say \emph{uṟuttu} again
\end{itemize}

\begin{tcolorbox}[breakable,colback=teal!4,colframe=teal!35!black,title={Before you move on}]
What does \ta{உறுத்து} mean? (\textquotedblleft{}To chafe; to irritate\textquotedblright{}.) Say it once more.
\end{tcolorbox}
